\documentclass[11pt]{article}
\usepackage[utf8]{inputenc}
\usepackage[T1]{fontenc}
\usepackage[margin=2.2cm]{geometry}
\usepackage{booktabs,adjustbox,amsmath,url}
\usepackage[hidelinks]{hyperref}
\newcommand{\DeltaSteels}{-0.3202}
\newcommand{\RelSteels}{-0.39}
\newcommand{\PNBSteels}{0.896}
\newcommand{\PBootSteels}{0.720}
\newcommand{\PHolmSteels}{1.000}
\newcommand{\FoldsSteels}{3}
\newcommand{\SteelsStack}{83.12}
\newcommand{\SteelsStackStd}{10.89}
\newcommand{\SteelsStackLocal}{82.80}
\newcommand{\SteelsStackLocalStd}{11.65}
\newcommand{\SteelsStackFm}{82.02}
\newcommand{\SteelsStackFmStd}{9.12}
\newcommand{\SteelsBestSingle}{83.41}
\newcommand{\SteelsBestSingleStd}{10.86}
\newcommand{\SteelsMeanEnsemble}{102.46}
\newcommand{\SteelsMeanEnsembleStd}{20.11}
\newcommand{\SteelsSingleLocal}{82.69}
\newcommand{\SteelsSingleLocalStd}{11.60}
\newcommand{\SteelsStackLocalVsBestSingleRel}{-0.74}
\newcommand{\SteelsStackLocalVsBestSingleP}{0.538}
\newcommand{\SteelsStackLocalVsBestSingleFolds}{3}
\newcommand{\SteelsStackVsBestSingleRel}{-0.35}
\newcommand{\SteelsStackVsBestSingleP}{0.660}
\newcommand{\SteelsStackVsBestSingleFolds}{2}
\newcommand{\SteelsSingleLocalVsBestSingleRel}{-0.86}
\newcommand{\SteelsSingleLocalVsBestSingleP}{0.425}
\newcommand{\SteelsSingleLocalVsBestSingleFolds}{3}
\newcommand{\SteelsStackLocalVsMeanEnsembleRel}{-19.19}
\newcommand{\SteelsStackLocalVsMeanEnsembleP}{$<$0.001}
\newcommand{\SteelsStackLocalVsMeanEnsembleFolds}{5}
\newcommand{\SteelsStackFmVsStackRel}{-1.32}
\newcommand{\SteelsStackFmVsStackP}{0.519}
\newcommand{\SteelsStackFmVsStackFolds}{2}
\newcommand{\SteelsStackFmVsStackLocalRel}{-0.94}
\newcommand{\SteelsStackFmVsStackLocalP}{0.583}
\newcommand{\SteelsStackFmVsStackLocalFolds}{3}
\newcommand{\SteelsStackLocalMagpieVsStackRel}{+0.50}
\newcommand{\SteelsStackLocalMagpieVsStackP}{0.036}
\newcommand{\SteelsStackLocalMagpieVsStackFolds}{0}
\newcommand{\SteelsStackLocalFracVsStackRel}{-0.39}
\newcommand{\SteelsStackLocalFracVsStackP}{0.721}
\newcommand{\SteelsStackLocalFracVsStackFolds}{3}
\newcommand{\SteelsNaiveOn}{5}
\newcommand{\SteelsFmOn}{5}
\newcommand{\SteelsPermNaiveOn}{5}
\newcommand{\SteelsPermFmOn}{5}
\newcommand{\SteelsFoldZeroStack}{91.07}
\newcommand{\SteelsFoldZeroLocal}{96.23}
\newcommand{\SteelsFoldZeroFm}{95.73}
\newcommand{\SteelsFoldZeroBest}{91.07}
\newcommand{\DeltaGap}{-0.00232}
\newcommand{\RelGap}{-0.71}
\newcommand{\PNBGap}{0.011}
\newcommand{\PBootGap}{0.001}
\newcommand{\PHolmGap}{0.005}
\newcommand{\FoldsGap}{5}
\newcommand{\GapStack}{0.3264}
\newcommand{\GapStackStd}{0.0139}
\newcommand{\GapStackLocal}{0.3240}
\newcommand{\GapStackLocalStd}{0.0140}
\newcommand{\GapStackFm}{0.3276}
\newcommand{\GapStackFmStd}{0.0149}
\newcommand{\GapBestSingle}{0.3259}
\newcommand{\GapBestSingleStd}{0.0130}
\newcommand{\GapMeanEnsemble}{0.4007}
\newcommand{\GapMeanEnsembleStd}{0.0242}
\newcommand{\GapSingleLocal}{0.3239}
\newcommand{\GapSingleLocalStd}{0.0133}
\newcommand{\GapStackLocalVsBestSingleRel}{-0.58}
\newcommand{\GapStackLocalVsBestSingleP}{0.071}
\newcommand{\GapStackLocalVsBestSingleFolds}{4}
\newcommand{\GapStackVsBestSingleRel}{+0.13}
\newcommand{\GapStackVsBestSingleP}{0.566}
\newcommand{\GapStackVsBestSingleFolds}{3}
\newcommand{\GapSingleLocalVsBestSingleRel}{-0.63}
\newcommand{\GapSingleLocalVsBestSingleP}{0.002}
\newcommand{\GapSingleLocalVsBestSingleFolds}{5}
\newcommand{\GapStackLocalVsMeanEnsembleRel}{-19.14}
\newcommand{\GapStackLocalVsMeanEnsembleP}{$<$0.001}
\newcommand{\GapStackLocalVsMeanEnsembleFolds}{5}
\newcommand{\GapStackFmVsStackRel}{+0.38}
\newcommand{\GapStackFmVsStackP}{0.390}
\newcommand{\GapStackFmVsStackFolds}{2}
\newcommand{\GapStackFmVsStackLocalRel}{+1.10}
\newcommand{\GapStackFmVsStackLocalP}{0.005}
\newcommand{\GapStackFmVsStackLocalFolds}{0}
\newcommand{\GapStackLocalMagpieVsStackRel}{-0.86}
\newcommand{\GapStackLocalMagpieVsStackP}{$<$0.001}
\newcommand{\GapStackLocalMagpieVsStackFolds}{5}
\newcommand{\GapStackLocalFracVsStackRel}{-0.88}
\newcommand{\GapStackLocalFracVsStackP}{0.001}
\newcommand{\GapStackLocalFracVsStackFolds}{5}
\newcommand{\GapNaiveOn}{5}
\newcommand{\GapFmOn}{5}
\newcommand{\GapPermNaiveOn}{5}
\newcommand{\GapPermFmOn}{5}
\newcommand{\GapFoldZeroStack}{0.3109}
\newcommand{\GapFoldZeroLocal}{0.3079}
\newcommand{\GapFoldZeroFm}{0.3107}
\newcommand{\GapFoldZeroBest}{0.3114}
\newcommand{\DeltaMetal}{-0.00000}
\newcommand{\RelMetal}{-0.00}
\newcommand{\PNBMetal}{0.541}
\newcommand{\PBootMetal}{0.899}
\newcommand{\PHolmMetal}{1.000}
\newcommand{\FoldsMetal}{0}
\newcommand{\MetalStack}{0.9825}
\newcommand{\MetalStackStd}{0.0017}
\newcommand{\MetalStackBA}{0.9374}
\newcommand{\MetalStackLocal}{0.9825}
\newcommand{\MetalStackLocalStd}{0.0018}
\newcommand{\MetalStackLocalBA}{0.9370}
\newcommand{\MetalStackFm}{0.9826}
\newcommand{\MetalStackFmStd}{0.0017}
\newcommand{\MetalStackFmBA}{0.9366}
\newcommand{\MetalBestSingle}{0.9819}
\newcommand{\MetalBestSingleStd}{0.0015}
\newcommand{\MetalBestSingleBA}{0.9378}
\newcommand{\MetalMeanEnsemble}{0.9811}
\newcommand{\MetalMeanEnsembleStd}{0.0027}
\newcommand{\MetalSingleLocal}{0.9819}
\newcommand{\MetalSingleLocalStd}{0.0015}
\newcommand{\MetalSingleLocalBA}{0.9378}
\newcommand{\MetalStackLocalVsBestSingleRel}{+0.06}
\newcommand{\MetalStackLocalVsBestSingleP}{0.301}
\newcommand{\MetalStackLocalVsBestSingleFolds}{4}
\newcommand{\MetalStackVsBestSingleRel}{+0.06}
\newcommand{\MetalStackVsBestSingleP}{0.275}
\newcommand{\MetalStackVsBestSingleFolds}{4}
\newcommand{\MetalSingleLocalVsBestSingleRel}{+0.00}
\newcommand{\MetalSingleLocalVsBestSingleP}{0.714}
\newcommand{\MetalSingleLocalVsBestSingleFolds}{2}
\newcommand{\MetalStackLocalVsMeanEnsembleRel}{+0.14}
\newcommand{\MetalStackLocalVsMeanEnsembleP}{0.008}
\newcommand{\MetalStackLocalVsMeanEnsembleFolds}{5}
\newcommand{\MetalStackFmVsStackRel}{+0.02}
\newcommand{\MetalStackFmVsStackP}{0.320}
\newcommand{\MetalStackFmVsStackFolds}{3}
\newcommand{\MetalStackFmVsStackLocalRel}{+0.02}
\newcommand{\MetalStackFmVsStackLocalP}{0.277}
\newcommand{\MetalStackFmVsStackLocalFolds}{3}
\newcommand{\MetalStackLocalMagpieVsStackRel}{+0.00}
\newcommand{\MetalStackLocalMagpieVsStackP}{1.000}
\newcommand{\MetalStackLocalMagpieVsStackFolds}{0}
\newcommand{\MetalStackLocalFracVsStackRel}{-0.00}
\newcommand{\MetalStackLocalFracVsStackP}{0.898}
\newcommand{\MetalStackLocalFracVsStackFolds}{0}
\newcommand{\MetalNaiveOn}{1}
\newcommand{\MetalFmOn}{5}
\newcommand{\MetalPermNaiveOn}{5}
\newcommand{\MetalPermFmOn}{5}
\newcommand{\MetalFoldZeroStack}{0.9850}
\newcommand{\MetalFoldZeroLocal}{0.9850}
\newcommand{\MetalFoldZeroFm}{0.9848}
\newcommand{\MetalFoldZeroBest}{0.9838}
\newcommand{\DeltaGlass}{+0.00000}
\newcommand{\RelGlass}{+0.00}
\newcommand{\PNBGlass}{1.000}
\newcommand{\PBootGlass}{1.000}
\newcommand{\PHolmGlass}{1.000}
\newcommand{\FoldsGlass}{0}
\newcommand{\GlassStack}{0.9567}
\newcommand{\GlassStackStd}{0.0091}
\newcommand{\GlassStackBA}{0.8841}
\newcommand{\GlassStackLocal}{0.9567}
\newcommand{\GlassStackLocalStd}{0.0091}
\newcommand{\GlassStackLocalBA}{0.8841}
\newcommand{\GlassStackFm}{0.9559}
\newcommand{\GlassStackFmStd}{0.0094}
\newcommand{\GlassStackFmBA}{0.8835}
\newcommand{\GlassBestSingle}{0.9548}
\newcommand{\GlassBestSingleStd}{0.0097}
\newcommand{\GlassBestSingleBA}{0.8827}
\newcommand{\GlassMeanEnsemble}{0.9458}
\newcommand{\GlassMeanEnsembleStd}{0.0116}
\newcommand{\GlassSingleLocal}{0.9548}
\newcommand{\GlassSingleLocalStd}{0.0097}
\newcommand{\GlassSingleLocalBA}{0.8827}
\newcommand{\GlassStackLocalVsBestSingleRel}{+0.20}
\newcommand{\GlassStackLocalVsBestSingleP}{0.002}
\newcommand{\GlassStackLocalVsBestSingleFolds}{5}
\newcommand{\GlassStackVsBestSingleRel}{+0.20}
\newcommand{\GlassStackVsBestSingleP}{0.002}
\newcommand{\GlassStackVsBestSingleFolds}{5}
\newcommand{\GlassSingleLocalVsBestSingleRel}{+0.00}
\newcommand{\GlassSingleLocalVsBestSingleP}{1.000}
\newcommand{\GlassSingleLocalVsBestSingleFolds}{0}
\newcommand{\GlassStackLocalVsMeanEnsembleRel}{+1.15}
\newcommand{\GlassStackLocalVsMeanEnsembleP}{$<$0.001}
\newcommand{\GlassStackLocalVsMeanEnsembleFolds}{5}
\newcommand{\GlassStackFmVsStackRel}{-0.09}
\newcommand{\GlassStackFmVsStackP}{0.153}
\newcommand{\GlassStackFmVsStackFolds}{1}
\newcommand{\GlassStackFmVsStackLocalRel}{-0.09}
\newcommand{\GlassStackFmVsStackLocalP}{0.149}
\newcommand{\GlassStackFmVsStackLocalFolds}{1}
\newcommand{\GlassStackLocalMagpieVsStackRel}{+0.00}
\newcommand{\GlassStackLocalMagpieVsStackP}{1.000}
\newcommand{\GlassStackLocalMagpieVsStackFolds}{0}
\newcommand{\GlassStackLocalFracVsStackRel}{+0.00}
\newcommand{\GlassStackLocalFracVsStackP}{1.000}
\newcommand{\GlassStackLocalFracVsStackFolds}{0}
\newcommand{\GlassNaiveOn}{0}
\newcommand{\GlassFmOn}{5}
\newcommand{\GlassPermNaiveOn}{5}
\newcommand{\GlassPermFmOn}{5}
\newcommand{\GlassFoldZeroStack}{0.9694}
\newcommand{\GlassFoldZeroLocal}{0.9694}
\newcommand{\GlassFoldZeroFm}{0.9699}
\newcommand{\GlassFoldZeroBest}{0.9690}
\newcommand{\NTasksImproved}{1}
\newcommand{\SteelsRank}{3}
\newcommand{\SteelsNEntries}{12}
\newcommand{\SteelsBoardBest}{TPOT-Mat}
\newcommand{\GapRank}{2}
\newcommand{\GapNEntries}{13}
\newcommand{\GapBoardBest}{Darwin}
\newcommand{\MetalAUCprob}{0.9825}
\newcommand{\MetalRank}{2}
\newcommand{\MetalNEntries}{8}
\newcommand{\MetalBoardBest}{Darwin}
\newcommand{\GlassAUCprob}{0.9567}
\newcommand{\GlassRank}{2}
\newcommand{\GlassNEntries}{8}
\newcommand{\GlassBoardBest}{MODNet (v0.1.12)}
\newcommand{\SteelsWallMin}{1}
\newcommand{\GapWallMin}{8}
\newcommand{\MetalWallMin}{4}
\newcommand{\GlassWallMin}{5}
\newcommand{\TotalCPUhWall}{1.5}
\newcommand{\SteelsScratchMin}{1}
\newcommand{\GapScratchMin}{16}
\newcommand{\MetalScratchMin}{8}
\newcommand{\GlassScratchMin}{44}
\newcommand{\SteelsPairShare}{44}
\newcommand{\GapPairShare}{33}
\newcommand{\MetalPairShare}{37}
\newcommand{\GlassPairShare}{84}
\newcommand{\SteelsLocalQOne}{+6.8}
\newcommand{\SteelsFmQOne}{-0.8}
\newcommand{\SteelsLocalQTwo}{-1.7}
\newcommand{\SteelsFmQTwo}{-2.8}
\newcommand{\SteelsLocalQThree}{-2.1}
\newcommand{\SteelsFmQThree}{+1.5}
\newcommand{\SteelsLocalQFour}{-1.5}
\newcommand{\SteelsFmQFour}{-2.2}
\newcommand{\GapLocalQOne}{+0.9}
\newcommand{\GapFmQOne}{+2.0}
\newcommand{\GapLocalQTwo}{-1.0}
\newcommand{\GapFmQTwo}{+0.5}
\newcommand{\GapLocalQThree}{-1.1}
\newcommand{\GapFmQThree}{+0.2}
\newcommand{\GapLocalQFour}{-0.9}
\newcommand{\GapFmQFour}{-0.3}
\newcommand{\MetalLocalQOne}{-0.1}
\newcommand{\MetalFmQOne}{+1.1}
\newcommand{\MetalLocalQTwo}{+0.0}
\newcommand{\MetalFmQTwo}{-1.0}
\newcommand{\MetalLocalQThree}{-0.2}
\newcommand{\MetalFmQThree}{+0.7}
\newcommand{\MetalLocalQFour}{-0.0}
\newcommand{\MetalFmQFour}{-0.6}
\newcommand{\GlassLocalQOne}{+0.0}
\newcommand{\GlassFmQOne}{+3.0}
\newcommand{\GlassLocalQTwo}{+0.0}
\newcommand{\GlassFmQTwo}{+0.8}
\newcommand{\GlassLocalQThree}{+0.0}
\newcommand{\GlassFmQThree}{-0.9}
\newcommand{\GlassLocalQFour}{+0.0}
\newcommand{\GlassFmQFour}{+0.8}

\title{Supplementary Material for ``A pre-registered test of cross-fitted stacking with a novelty-damped local residual correction on the Matbench composition tasks''}
\author{Bhaskar Gurram (Independent Researcher)}
\date{}
\begin{document}
\maketitle

\section*{S1. Pre-registration and amendment}
The analysis plan (\texttt{ANALYSIS\_PLAN.md}) and SHA-256 hashes of all code were committed before any official test fold was scored.
\begin{itemize}
\item \textbf{Original freeze:} \texttt{results/freeze\_hashes\_original\_plan.txt}. Under it, the design was fixed on inner CV of fold~0's training rows.
\item \textbf{Amendment~1:} \texttt{results/freeze\_hashes.txt}, plan Section~8. Written before any scoring. It moved the choice among the explored design configurations inside every official training fold, because fold~0's training rows contain the test rows of folds 1--4.
\item \textbf{Superseded run:} a first run under the original plan produced test predictions that were never scored. They are archived unscored in the repository.
\item \textbf{Deviations:} all are listed in plan Section~7. They are a one-line fix in the scoring script (iterating over each fold's own arms), scheduling-only changes, and reporting scripts.
\end{itemize}

\section*{S1b. Amendment 2}
Amendment~2 (plan Sections 9--11, \texttt{results/freeze\_hashes\_amendment2.txt}) was written after the Amendment-1 test results had been seen. It changed the expert pool, the features (including a disclosed DFT transfer block) and the correction gate. It was frozen before any Amendment-2 test prediction, and the test folds were scored once under it. The design evidence (fold~0 training rows only) is in \texttt{results/design\_a2/DESIGN\_LOG\_A2.md}.
The leaders' per-sample predictions used for the paired comparison come from the public \texttt{matbench} repository (\texttt{benchmarks/matbench\_v0.1\_\{TPOT,darwin,modnet\_v0.1.12\}/results.json.gz}).

\section*{S2. Per-fold design selection (Amendment 1)}
Table~\ref{tab:ds} lists the candidate selected inside every official training fold, with the cross-fitted OOF score of the uncorrected global stack for each candidate (MAE for regression; ROC-AUC for classification).
Regression candidates:
\begin{itemize}
\item A: base features, L2 stack;
\item B: base features, L1 stack, plus an L1-objective LightGBM expert;
\item D: extra features, L1 stack;
\item C: extra features, L1 stack, plus the L1-LightGBM expert.
\end{itemize}
Classification candidates: A (base features) and C (extra features).

\begin{table}[h]
\centering
\caption{Per-fold design selection (inside each official training fold).}
\label{tab:ds}
\begin{adjustbox}{max width=\textwidth}
\begin{tabular}{llll}
\toprule
Task & Fold & Selected design & Candidate stack OOF scores \\
\midrule
steels & 0 & B\_base\_l1 & A=84.3708, B=83.1777, D=85.4401, C=85.4401 \\
steels & 1 & D\_extra\_l1 & A=90.1208, B=89.8331, D=89.2622, C=89.2622 \\
steels & 2 & A\_base\_l2 & A=92.3679, B=94.5285, D=93.8040, C=93.8040 \\
steels & 3 & C\_extra\_l1\_lgbml1 & A=92.6136, B=92.0424, D=91.4351, C=91.4351 \\
steels & 4 & B\_base\_l1 & A=88.8951, B=88.8950, D=89.2953, C=89.2953 \\
expt\_gap & 0 & D\_extra\_l1 & A=0.4159, B=0.4121, D=0.3408, C=0.3408 \\
expt\_gap & 1 & D\_extra\_l1 & A=0.4192, B=0.4180, D=0.3408, C=0.3409 \\
expt\_gap & 2 & D\_extra\_l1 & A=0.4067, B=0.4021, D=0.3411, C=0.3411 \\
expt\_gap & 3 & D\_extra\_l1 & A=0.4175, B=0.4185, D=0.3451, C=0.3454 \\
expt\_gap & 4 & D\_extra\_l1 & A=0.4115, B=0.4062, D=0.3385, C=0.3385 \\
expt\_is\_metal & 0 & C\_extra & A=0.9739, C=0.9819 \\
expt\_is\_metal & 1 & C\_extra & A=0.9733, C=0.9808 \\
expt\_is\_metal & 2 & C\_extra & A=0.9746, C=0.9813 \\
expt\_is\_metal & 3 & C\_extra & A=0.9725, C=0.9795 \\
expt\_is\_metal & 4 & C\_extra & A=0.9743, C=0.9812 \\
glass & 0 & C\_extra & A=0.9458, C=0.9461 \\
glass & 1 & A\_base & A=0.9521, C=0.9520 \\
glass & 2 & C\_extra & A=0.9515, C=0.9523 \\
glass & 3 & C\_extra & A=0.9508, C=0.9508 \\
glass & 4 & C\_extra & A=0.9449, C=0.9453 \\
\bottomrule
\end{tabular}

\end{adjustbox}
\end{table}

\section*{S3. Ideation and prior-art process}
Literature search, prior-art checking and idea generation used the ResearchStudio Idea skills (paper-search, scoop-check, idea-spark), driven by an AI assistant; all records are in the repository (\texttt{research/}, \texttt{ideaspark\_run/}).
\begin{itemize}
\item \textbf{Sources.} Retrieval used arXiv, OpenAlex, Semantic Scholar and Crossref. DBLP failed to return parseable responses, and OpenReview returned no records (rate-limited and off-venue).
\item \textbf{Prior-art verdict: medium overlap.}
  \begin{itemize}
  \item No retrieved work combines cross-fitted heterogeneous stacking, a novelty-damped local OOF correction and the Matbench composition tasks.
  \item Every component has ancestors: stacking, regression kriging and $\Delta$-learning.
  \item The closest external works are a stack with a global kernel-ridge residual correction in operator learning, MODNet, and a stacked band-gap ensemble.
  \end{itemize}
\item \textbf{Idea card.} idea-spark proposed the fold-matched correction. Its coherence check verified the decomposition naive $=$ analogue $+$ absorption on synthetic instances, and its independent post-revision review passed.
\item \textbf{Run outcome.} The run ended in a \texttt{needs\_work} state because of an internal inconsistency in the tool's novelty re-audit gate. We therefore treated the proposal as unvalidated and pre-registered it as a secondary arm.
\end{itemize}

\section*{S4. Design-phase log}
All configurations explored on fold~0's training rows (five pseudo-folds) are listed in \texttt{results/design/DESIGN\_LOG.md}.
These numbers informed only which candidates exist, not the test-fold results.

\section*{S5. Leaderboard snapshot and label-based classification scores}
The leaderboard values used (\texttt{results/leaderboard\_snapshot.json}) were taken from the per-task Matbench pages on 18 October 2026.

For every listed classification entry, the reported mean ``ROC-AUC'' equals the mean balanced accuracy to four decimals. This is what \texttt{matbench} computes when boolean labels are recorded: ROC-AUC computed from hard labels equals balanced accuracy.

We provide two unsubmitted result files: one with labels at the inner-CV threshold (comparable with the board) and one with probabilities.

\end{document}
